% =============================================================
%  upLaTeX 統合ショーケース（jsarticle + dvipdfmx）
%  学会テンプレ系の組版経路（upLaTeX -> dvipdfmx）で主要パッケージをまとめて使用。
%  dvipdfmx ドライバはクラスオプションで全パッケージに伝播させる。
%  ビルド: latexmk showcase.tex  →  out/showcase.pdf
% =============================================================
\documentclass[uplatex,dvipdfmx,a4paper,11pt]{jsarticle}

\usepackage{amsmath,amssymb}
\usepackage{graphicx}
\usepackage{booktabs}
\usepackage{array}
\usepackage{xcolor}
\usepackage{enumitem}
\usepackage[margin=25mm]{geometry}
\usepackage{caption}
\usepackage{listings}
\usepackage{siunitx}
\usepackage[hidelinks]{hyperref}
\usepackage{cleveref}

\lstset{basicstyle=\ttfamily\small, frame=single}

\title{upLaTeX 統合ショーケース}
\author{latex-ja-template テスト}
\date{\today}

\begin{document}
\maketitle

\section{はじめに}
このドキュメントは、学会テンプレートと同じ upLaTeX $\to$ dvipdfmx の経路で、
主要パッケージをまとめて使用し PDF まで生成できることを確認します。
色付きテキスト \textcolor{red}{赤}・\colorbox{yellow}{ハイライト} も使えます。

\section{数式}
\begin{align}
  E &= mc^2, \label{eq:e} \\
  \int_{-\infty}^{\infty} e^{-x^2}\,dx &= \sqrt{\pi}. \label{eq:gauss}
\end{align}
\cref{eq:e,eq:gauss} を cleveref で参照しています。集合 $\mathbb{R}$、不等号 $\leq$ も表示できます。

\section{箇条書き}
\begin{enumerate}[label=(\arabic*), noitemsep]
  \item enumitem によるラベル変更。
  \item 項目間隔の調整。
\end{enumerate}

\section{表}
\begin{table}[h]
  \centering
  \caption{booktabs と array による表。}
  \label{tab:show}
  \begin{tabular}{>{\bfseries}l c c}
    \toprule
    手法 & 精度 & 速度 \\
    \midrule
    手法A & 0.91 & 速い \\
    手法B & 0.95 & 普通 \\
    \bottomrule
  \end{tabular}
\end{table}
\cref{tab:show} を参照。

\section{図}
\begin{figure}[h]
  \centering
  \scalebox{1.0}{\color{blue}\rule{4cm}{2cm}}
  \caption{graphicx の \texttt{scalebox} と \texttt{rule} で生成した図。}
  \label{fig:show}
\end{figure}

\section{単位}
速度 \qty{9.8}{\meter\per\second\squared}、個数 \num{12345.678}。

\section{コード}
\begin{lstlisting}[language=Python]
def hello():
    print("Hello, LaTeX")
\end{lstlisting}

\section{リンク}
詳細は \href{https://example.com}{example.com} を参照してください。

\end{document}
